% Source: https://artofproblemsolving.com/community/c3h1748211_tough_algebra_problem
% Stable identity: aops-topic-1748211-q01
Write $(x-y)(x+y)=2000^2=2^8\cdot5^6$. The factors $x-y$ and $x+y$ have the same parity (their sum is $2x$); their product is even, so both are even. Put $x-y=2m$ and $x+y=2n$; then
\[
mn=2^6\cdot5^6=10^6,
\]
and conversely every pair of integers $(m,n)$ with $mn=10^6$ gives the lattice point $(x,y)=(m+n,\,n-m)$, different pairs giving different points (for example $m=250$, $n=4000$ gives $(4250,3750)$, and $m=500$, $n=2000$ gives $(2500,1500)$).

\begin{asy}
usepackage("amsmath");
size(6.5cm);
defaultpen(fontsize(10pt));
import graph;
pair hp(real t){return (2*cosh(t),2*sinh(t));} pair hm(real t){return (-2*cosh(t),2*sinh(t));}
draw((-4.2,0)--(4.4,0),gray,Arrow(5));draw((0,-3.4)--(0,3.6),gray,Arrow(5));
label("$x$",(4.4,0),E,gray);label("$y$",(0,3.6),N,gray);
draw((-3.3,-3.3)--(3.3,3.3),gray+dashed);draw((-3.3,3.3)--(3.3,-3.3),gray+dashed);
draw(graph(hp,-1.15,1.15),blue+linewidth(0.9));draw(graph(hm,-1.15,1.15),blue+linewidth(0.9));
pair[] P={(2,0),(-2,0),(2.5,1.5),(2.5,-1.5),(-2.5,1.5),(-2.5,-1.5)};
for(pair z:P) dot(z);
label("$2000$",(1.25,-0.3),fontsize(8pt));label("$-2000$",(-1.2,-0.3),fontsize(8pt));
label("$(2500,1500)$",(2.5,1.5),E);label("$(2500,-1500)$",(2.5,-1.5),E);
label("\mbox{}",(min(currentpicture,user=true)+max(currentpicture,user=true))/2,black); // keeps the text after the figure black when LaTeX typesets the labels inline
\end{asy}

The number $10^6$ has $(6+1)(6+1)=49$ positive divisors, and $m$ may be any positive or negative divisor (then $n=\dfrac{10^6}{m}$), so there are $2\cdot49=\boxed{98}$ lattice points.
